\section{Rank-tight compatibility and simplex geometry}\label{sec:simplex50}
Throughout this section $D\ge2$, $X=\{\pm e_1,\ldots,\pm e_D\}$,
all coordinate queries are available, and the orthogonal alphabet contains
$I$. The signal satisfies $0<\rho<1$. The conclusions concern exact
realization unless a positive tolerance is explicitly specified. Put
\[
 C_\rho=\rho\conv X,\qquad
 Q_t=\{z\in\R^D:|\langle e_j,U_vz\rangle|\le1
       \text{ for }|v|=N-t,\ 1\le j\le D\}.
\]
Identity padding gives $Q_t\subset[-1,1]^D$. Define the injective affine
map $J_t:Q_t\to\mathcal C_t$ by
\[
 J_t(z)_{v,j,b}=(1+b\langle e_j,U_vz\rangle)/2.
\]
The spaces $\mathcal C_t$, the restrictions $L_{t,a}$ and the normalized
Hankel arrays are those of Section~\ref{sec:hankel}. In particular,
$L_{t,a}J_t(z)=J_{t+1}(U_az)$.

\begin{theorem}[Rank-tight simplex equivalence]\label{thm:simplex50}
An exact clocked realization with at most $D+1$ labels at every command
cut exists if and only if there are full-dimensional simplices
$S_t\subset Q_t$, each with $D+1$ vertices, such that
\begin{equation}\label{eq:simplex50}
 C_\rho\subset S_0,\qquad U_aS_t\subset S_{t+1}
       \quad(0\le t<N,\ a\in\calA).
\end{equation}
More generally, at every cut with $D+1$ labels in an arbitrary exact
machine, every label's full future-response row is $J_t(z)$ for a
unique $z\in Q_t$. If $s<t$ are two such cuts, their simplices satisfy
$U_wS_s\subset S_t$ for every word of length $t-s$, irrespective of
all intervening widths.
\end{theorem}
\begin{proof}
Every $H_t$ has rank $D+1$. Its rows are affine functions of
$\rho U_ux$, so its rank is at most $D+1$. Identity prefixes and suffixes,
the signed coordinate seeds, and the coordinate queries give rank
$D+1$. At a cut of this width the machine factorization $H_t=E_tB_t$
has $E_t$ of full column rank and $B_t$ of full row rank. Hence
$\operatorname{rowspan}B_t=\operatorname{rowspan}H_t$. Write a row of
$B_t$ as a linear combination of rows of $H_t$. Normalization of any
binary column pair forces the sum of its coefficients to be one.
It is consequently an affine combination and equals $J_t(z)$.
Identity suffixes give uniqueness of $z$; legality gives $z\in Q_t$.
The $D+1$ rows are linearly independent, so their $z$'s are affinely
independent. The initialization contains $C_\rho$ in their convex hull.
Composing the actual transition rows between two narrow cuts gives
\[
 L_wB_s=T_wB_t.
\]
Injectivity of $J_t$ gives the asserted containment for every generator,
including generators unreachable from a particular seed. No observable
projection of a higher-width row has been used.

Conversely express every $\rho x$ in the vertices of $S_0$, and every
$U_av$ in the vertices of $S_{t+1}$, by barycentric coordinates. They
are nonnegative and sum to one, hence are legal initialization and
transition rows. Decode coordinate $j$ by the $j$th coordinates of
the terminal vertices. These lie in $[-1,1]$. Induction gives the
conditional mean $\rho U_wx$ and all required binary laws. The extra
held-answer cut has two labels. Rank already excludes fewer than
$D+1$ labels at any command cut.
\end{proof}

At a rank-tight cut, full column rank of $E_t$ also excludes a label unreachable from every prefix: such a label would give a zero column. The rank hypothesis is decisive. For a larger factorization,
$\operatorname{rowspan}B_t$ may strictly contain that of $H_t$, and the
preceding elimination of hidden future directions is invalid. The
underlying cone/factorization principle is classical \cite{BF04}; the
rank-tight conclusion identifies precisely when the observed-space
simplex description is necessary as well as sufficient in this
nonuniform controlled problem.

\subsection{A finite occupation certificate}
Let
\[
 c_D=2^{-D}\binom{2D}{D}>1\quad(D\ge2).
\]
We record the elementary simplex case of the difference-body identity;
this is classical convex geometry, not a new volume inequality.

\begin{lemma}[Difference body of a simplex]\label{lem:difference50}
For a full-dimensional simplex $S$,
$\operatorname{vol}(S-S)=\binom{2D}{D}\operatorname{vol}(S)$.
If a convex body $T$ contains both $S$ and $-S$, then
$\operatorname{vol}(T)\ge c_D\operatorname{vol}(S)$.
\end{lemma}
\begin{proof}
Affine covariance reduces the identity to
$\Delta=\{x\ge0:\sum x_i\le1\}$. Its difference body is given by
$\sum x_i^+\le1$ and $\sum x_i^-\le1$. An orthant with $k$ positive
coordinates contributes $1/(k!(D-k)!)$. Summing its $\binom{D}{k}$
occurrences and using $\sum_k\binom{D}{k}^2=\binom{2D}{D}$ proves the
identity, since $\operatorname{vol}(\Delta)=1/D!$. The containment
$(S-S)/2\subset\conv(S\cup(-S))\subset T$ proves the inequality.
\end{proof}

\begin{theorem}[Rank-tight occupation budget]\label{thm:occupation50}
Suppose $I,-I\in\calA$. In every exact realization, let $m$ be the
number of command cuts whose available width is $D+1$, including cuts $0$ and $N$ when they have that width. If $m\ge1$, then
\begin{equation}\label{eq:occupation50}
 c_D^{\,m-1}\le D!\rho^{-D}.
\end{equation}
All other cuts may have arbitrary finite widths. In particular, an
all-$(D+1)$ profile at horizon $N$ requires
$c_D^N\le D!\rho^{-D}$.
\end{theorem}
\begin{proof}
Order the narrow cuts. Between consecutive ones the identity word and
a word with exactly one $-I$ and all other letters equal to $I$ exist. Theorem~\ref{thm:simplex50} therefore
forces the next simplex to contain both signs of the preceding one.
Lemma~\ref{lem:difference50} multiplies volume by at least $c_D$.
At the first narrow cut identity prefixes give $C_\rho\subset S$,
so its volume is at least $(2\rho)^D/D!$. At the last narrow cut
identity suffixes give $S\subset[-1,1]^D$, of volume $2^D$.
Comparing these bounds proves the result. Intermediate registers and
their private randomness have disappeared only through their actual
composed stochastic rows, not through an uncharged history register.
\end{proof}

\begin{corollary}[An exact four-state frontier]\label{cor:four50}
In dimension two let $\calA$ be any alphabet of signed permutation
matrices containing $I,-I$. For every $0<\rho<1$ and every horizon with
\begin{equation}\label{eq:four50}
 (3/2)^N>2\rho^{-2},
\end{equation}
one has $W_{N,0}=4$. For $\rho=1/10$, all $N\ge14$ satisfy this
condition. The bound $14$ is sufficient; it is not asserted to be the
first horizon requiring four labels. For each horizon satisfying
\eqref{eq:four50}, four labels remain optimal throughout some interval
$0\le\varepsilon<\eta_N(\calA,\rho)$, where $\eta_N(\calA,\rho)>0$.
\end{corollary}
\begin{proof}
Three labels at every cut violate Theorem~\ref{thm:occupation50}; fewer
labels are excluded by Hankel rank. Four labels suffice by storing
$x\in\{\pm e_1,\pm e_2\}$ exactly, permuting that label under each
command, and decoding by $\rho x_j$. This is a common-row construction
with no horizon-dependent transition. The integer comparison
$3^{14}>200\,2^{14}$ proves the numerical assertion. At fixed $N$ the
space of machines with at most three labels is represented by padding
to exactly three and is compact. The continuous worst-row error has a
strictly positive minimum $\eta_N(\calA,\rho)$, since zero is excluded. This gives
the positive-tolerance assertion without assigning an unproved
numerical value to $\eta_N(\calA,\rho)$.
\end{proof}

For example, adjoining the swap matrix and
$\operatorname{diag}(1,-1)$ gives genuinely noncommuting commands; their
products in opposite orders differ. This is a class-wide causal
occupation obstruction beyond width two, followed by an attaining
machine, not a test of supplied transition rows. The exact occupation
constant need not be sharp.

\subsection{From arbitrary horizons to a single permutation action}
\begin{theorem}[All-horizon simplex rigidity]\label{thm:allhorizon50}
Under the hypotheses at the start of this section, the following are
equivalent: every finite horizon admits an exact realization of peak
$D+1$; there is a single full-dimensional simplex
$C_\rho\subset S\subset[-1,1]^D$ with $U_aS=S$ for all letters.
In the latter case one realization works for all horizons and all
command updates are permutations of its $D+1$ vertices. The group
generated by the commands is finite and embeds in the symmetric group
on $D+1$ letters.
\end{theorem}
\begin{proof}
Take simplex chains at successively increasing horizons. At every cut
identity prefixes and suffixes place their vertices in the compact
cube and their convex hulls around $C_\rho$. A diagonal subsequence of
vertex tuples produces an infinite chain $S_0,S_1,\ldots$ satisfying
\eqref{eq:simplex50}. This first extraction runs over horizons. The enclosed crosspolytope prevents degeneracy.
Since $I$ is a letter, this chain is nested. Its Hausdorff limit
$S_\infty=\overline{\bigcup_t S_t}$ is again a simplex with at most
$D+1$ vertices: choose a convergent subsequence of the bounded vertex
tuples along this one nested chain. This is the second extraction. It remains full-dimensional. Continuity gives
$U_aS_\infty\subset S_\infty$. Orthogonality preserves volume, so this
containment is equality. Indeed, a proper inclusion of full-dimensional
compact convex bodies strictly decreases volume. Each $U_a$ therefore
permutes the vertices. An affine map fixing $D+1$ independent vertices
is the identity, which proves faithfulness and finiteness.
Conversely barycentric initialization, the vertex permutations, and
coordinate decoders on $S$ give one exact machine at every horizon.
All its suffix responses are legal because every group element
preserves $S\subset[-1,1]^D$.
\end{proof}

In the plane, for $\rho\le1/2$, this yields a complete criterion for
three-state feasibility at all horizons: the generated group must be
contained in the symmetry group of some equilateral triangle. Necessity
follows by classifying the orthogonal symmetries of a triangle: a
nontrivial rotation has order three and forces an equilateral triangle;
without such a rotation the group has at most one reflection. For
sufficiency a regular triangle of circumradius one and the appropriate
orientation contains the disk of radius $1/2$. In particular, order-three
rotations together with a compatible reflection give a noncommutative
three-state class. A single reflection is not a half-turn: the latter cannot preserve a triangle. Two distinct reflections in the triangular group have product of order three. The arbitrary-width counterpart is Theorem~\ref{thm:finitegroup51}.

\section{An exact separation of distortion and enclosure}\label{sec:separation50}
We now link the rank-tight theorem to the inherited word profiles.
The obstruction is not a better choice of a packet law: it is information
that every one-direction packet law omits.

\begin{theorem}[Strict distortion--dilation separation]\label{thm:separation50}
For $\calA=\{I,-I\}$ in dimension $D\ge2$,
\[
 \Gamma_{\calA}(k,B)=0\quad(k\ge2,\ B\ge1),\qquad
 \mathfrak r_{\calA}(k)=0\quad(k\ge2).
\]
Here $a$ is an inradius parameter, not a command letter. Nevertheless, for every $0<a\le1/D$ the common enclosure profile is
\begin{equation}\label{eq:lambda50}
 \lambda_{\calA}(D+1,a)=D.
\end{equation}
Thus the inequality in Theorem~\ref{thm:rate44} is strictly unequal:
$0<\log D$. At rank-tight cuts the zero packet charges coexist with
the positive occupation charge $\log c_D$ from
Theorem~\ref{thm:occupation50}.
\end{theorem}
\begin{proof}
For each unit vector $u$, every word sends $u$ to one of $u,-u$.
Taking these two vectors as centers makes every scalar product maximum
one, for every word law. The definitions of $D_k$, $\Gamma$ and
$\mathfrak r$ give their asserted values.
A full-dimensional polytope with at most $D+1$ vertices is a simplex.
Let its vertices be $v_i$, and write the barycentric coordinates of
zero as $c_i>0$, $\sum c_i=1$. If $-S\subset\lambda S$, uniqueness of
barycentric coordinates for $-v_i/\lambda$ forces
\[
 (1+\lambda^{-1})c_i-\lambda^{-1}\ge0,
 \quad\text{hence}\quad c_i\ge(1+\lambda)^{-1}.
\]
Summing over $D+1$ vertices yields $\lambda\ge D$. A regular simplex
centered at zero with circumradius one has inradius $1/D$ and obeys
$-S\subset D S$ by the same coordinates $c_i=1/(D+1)$.
It is feasible for every stated $a$, proving equality. The occupation
comparison follows from the previous theorem, not from an invalid
reversal of the packet bound.
\end{proof}

This exact separation preserves rather than replaces the arithmetic
results below. They use situations where distortion is positive and
can match enclosure scales. Here the two mechanisms are provably
inequivalent, and rank-tight Hankel rigidity supplies the missing
finite-dimensional geometric obstruction. The value $D$ is the
classical simplex asymmetry calculation; its role here is the exact
comparison with the optimized controlled-word profile.
